\documentclass[../../../main.tex]{subfiles}
\begin{document}
\subsection{Fidelity}
The fidelity of quantum channel is
\begin{equation*}
  \braket{F }  = \frac{1 }{2 } + \frac{1 }{3} \mathrm{Re} (f_r) +\frac{1 }{6 }|f_r|^2
\end{equation*}
Classical fidelity is
\begin{equation*}
  \braket{F } = \frac{1 }{2 } \sum_{k}a_k + \frac{1 }{6 }\sum_{k } \mathbf{b }_k \cdot \mathbf{r}_k
\end{equation*}
This is in general unoptimized, with the maximum of $\braket{F} = 2/3$
Another relevant quantum state transition amplitude
\begin{equation*}
  f_j(t) = \bra{\mathbf{j}} e ^{-iHt/\hbar }\ket{\mathbf{s}}
\end{equation*}

Quantum fidelity measured by averaging the output reduced density state with respect to input state $F=\int d\Omega \braket{\psi | \rho_\text{out}|\psi}$, with 
\begin{equation*}
  \rho_\text{out}
  = \begin{bmatrix}
    1- |\beta|^2 |f_r|^2 & \alpha\beta^*  f_r^* \\
    \alpha\beta  f_r     & |\beta|^2 |f_r|^2
  \end{bmatrix}
\end{equation*}
This state has the purity of 
\begin{equation*}
    \mathrm{Tr} [\rho_\text{out }^2] = 1+ 2|\beta|^{4}|f_{r}|^{2}(|f_r| ^2 -1)  
\end{equation*}

\subsubsection{Derivation: quantum fidelity}
Assume all down ground state
\begin{equation*}
  \ket{\mathbf{0 } } = \bigotimes_{i=1}^N \ket{0_i} = \ket{\downarrow_1\ldots \downarrow_N}
\end{equation*}
and state $\ket{\mathbf{j}}$ such that $j$-th site has been excited, with
\begin{equation*}
  \mathbf{j}=\mathbf{1},\dots,\mathbf{s},\dots,\mathbf{r},\dots,\mathbf{N}
\end{equation*}
Then the sender Alice places a spin in the unknown state
\begin{equation*}
  \ket{\psi_\text{in }} = \cos \frac{\theta }{2 }\ket{0} + e ^{i\phi} \sin \frac{\theta }{2 } \ket{1} = \alpha \ket{0} + \beta \ket{1}
\end{equation*}
at the $s$-th site in the spin chain.
We can describe the state of the whole chain at this instant
\begin{equation*}
  \ket{\Psi(0)} = \cos \frac{\theta }{2 }\ket{\mathbf{0 }} + e ^{i\phi} \sin \frac{\theta }{2 } \ket{\mathbf{s}}
\end{equation*}
Do note the difference of Alice encoded state $\ket{\psi_\text{in}}$ with the whole spin chain state $\ket{\Psi(t)}$, also $\ket{\mathbf{s}}$ denote the state at which $s$-th site, which is Alice's, is excited.
As $[H,\sigma_3]=0$, the total 3rd component is conserved and state $\ket{\mathbf{s}}$ evolve into $\ket{\mathbf{j}}$
\begin{align*}
  \ket{\Psi(t)}
   & =  \cos \frac{\theta }{2 } e ^{-iHt/\hbar } \ket{\mathbf{0 }} + e ^{i\phi} \sin \frac{\theta }{2 } e ^{-iHt/\hbar }\ket{\mathbf{s}}                                    \\
   & =  \cos \frac{\theta }{2 }  \ket{\mathbf{0 }} + e ^{i\phi} \sin \frac{\theta }{2 } \sum_{j=1 } ^N \ket{\mathbf{j } } \bra{\mathbf{j}} e ^{-iHt/\hbar }\ket{\mathbf{s}} \\
  \ket{\Psi(t)}
   & = \alpha \ket{\mathbf{0 }} + \beta \sum_{j=1 } ^N f_j(t) \ket{\mathbf{j}}
\end{align*}
with the ground state energy set as zero, $\hbar = 1$ and transition amplitude $f_j(t) = \bra{\mathbf{j}} e ^{-iHt/\hbar }\ket{\mathbf{s}}$.

The density matrix in this case
\begin{align*}
  \rho(t) & =
  \left(\alpha \ket{\mathbf{0 }} + \beta \sum_{j=1 } ^N f_j(t) \ket{\mathbf{\mathbf{j}}} \right)
  \left(\alpha \bra{\mathbf{0 }} + \beta \sum_{k=1 } ^N f_k(t) ^* \bra{\mathbf{\mathbf{k}}} \right)                                                                                                                                                    \\
  \rho(t) & = \alpha^2   \ket{\mathbf{0}}\bra{\mathbf{0}} + \alpha \beta^* \sum_k f_k^*  \ket{\mathbf{0}}\bra{\mathbf{k}} + \alpha \beta \sum_j f_j \ket{\mathbf{j}}\bra{\mathbf{0}} + |\beta|^2 \sum_{j,k} f_j f_k^* \ket{\mathbf{j}}\bra{\mathbf{k}}
\end{align*}
since $\alpha$ is real.
To obtain the sender Bob state, we use the definition of partial trace
\begin{equation*}
  \rho_{\text{out}} (t) = \mathrm{Tr} _{A'} [\rho(t)]
  = \sum_{\boldsymbol{\ell}} \left( \mathbf{1 }_\mathbf{r}\otimes\bra{\boldsymbol{\ell}}_{A'} \right) \rho \left( \mathbf{1 }_\mathbf{r}\otimes\ket{\boldsymbol{\ell}}_{A'} \right)
\end{equation*}
with $A' = \{\mathbf{j}\}\setminus \mathbf{r}$.
The basis can be factored in terms of these subspaces
\begin{equation*}
  \ket{\mathbf{0} } = \ket{0}_\mathbf{r} \otimes \ket{\mathbf{0} }_{A'}
  \qquad
  \ket{\mathbf{j} } =
  \begin{cases}
    \ket{1}_\mathbf{r} \otimes \ket{\mathbf{0} }_{A'} & \mathbf{j } = \mathbf{r}     \\
    \ket{0}_\mathbf{r} \otimes \ket{\mathbf{j} }_{A'} & \mathbf{j } \neq  \mathbf{r} \\
  \end{cases}
\end{equation*}
Unlike $\ket{\mathbf{j}}$, the $\ket{\boldsymbol{\ell}}_{A'}$ spans the full Hilbert space of $A'$, not just the one-excitation sector
\begin{equation*}
  \ket{\boldsymbol{\ell}}_{A'} = \ket{\ell,\dots,\ell_{r-1},\ell_{r+1},\dots,\ell_{N}}, \qquad\ell_i \in \{0,1\}
\end{equation*}
The partial trace can be performed term by term, starting with the vacuum term
\begin{align*}
  \mathrm{Tr} _{A'} \ket{\mathbf{0}}\bra{\mathbf{0}}
   & = \sum_{\boldsymbol{\ell}} \left( \mathbf{1 }_\mathbf{r}\otimes\bra{\boldsymbol{\ell}}_{A'} \right)\ket{\mathbf{0}}\bra{\mathbf{0}} \left( \mathbf{1 }_\mathbf{r}\otimes\ket{\boldsymbol{\ell}}_{A'} \right) \\
  \mathrm{Tr} _{A'} \ket{\mathbf{0}}\bra{\mathbf{0}}
   & = \sum_{\boldsymbol{\ell}} \ket{0}_\mathbf{r}  \braket{\boldsymbol{\ell}| \mathbf{0}}_{A'} \braket{\mathbf{0}| \boldsymbol{\ell}}_{A'}
  = \ket{0}_\mathbf{r} \bra{0}_\mathbf{r}
\end{align*}
For the cross term
\begin{equation*}
  \mathrm{Tr} _{A'} \ket{\mathbf{0}}\bra{\mathbf{j}}
  = \sum_{\boldsymbol{\ell}} \left( \mathbf{1 }_\mathbf{r}\otimes\bra{\boldsymbol{\ell}}_{A'} \right)\ket{\mathbf{0}}\bra{\mathbf{j}} \left( \mathbf{1 }_\mathbf{r}\otimes\ket{\boldsymbol{\ell}}_{A'} \right)
\end{equation*}
there are two possibilities, at $\mathbf{j} = \mathbf{r}$
\begin{align*}
  \mathrm{Tr} _{A'} \ket{\mathbf{0}}\bra{\mathbf{j}}
   & =  \sum_{\boldsymbol{\ell}} \left( \mathbf{1 }_\mathbf{r} \bra{\boldsymbol{\ell}}_{A'} \right)\left(\ket{0}_\mathbf{r} \ket{\mathbf{0}}_{A'} \right) \left( \bra{{1}}_\mathbf{r}  \ket{\mathbf{0}}_{A'}\right)  \left( \mathbf{1 }_\mathbf{r} \ket{\boldsymbol{\ell}}_{A'} \right) \\
   & = \sum_{\boldsymbol{\ell}}  \ket{0}_\mathbf{r} \braket{\boldsymbol{\ell  }  |\mathbf{0} }_{A'} \bra{1}_\mathbf{r} \braket{\mathbf{0 } | \boldsymbol{\ell }}_{A'}
  = \ket{0}_\mathbf{r} \bra{1}_\mathbf{r}
\end{align*}
and at $\mathbf{j} \neq \mathbf{r}$
\begin{align*}
  \mathrm{Tr} _{A'} \ket{\mathbf{0}}\bra{\mathbf{j}}
   & =  \sum_{\boldsymbol{\ell}} \left( \mathbf{1 }_\mathbf{r} \bra{\boldsymbol{\ell}}_{A'} \right)\left(\ket{0}_\mathbf{r} \ket{\mathbf{0}}_{A'} \right) \left( \bra{{0}}_\mathbf{r}  {\mathbf{j}}_{A'}\right)  \left( \mathbf{1 }_\mathbf{r} \ket{\boldsymbol{\ell}}_{A'} \right) \\
   & = \sum_{\boldsymbol{\ell}}  \ket{0}_\mathbf{r} \braket{\boldsymbol{\ell  }  |\mathbf{0} }_{A'} \bra{1}_\mathbf{r} \braket{\mathbf{j } | \boldsymbol{\ell }}_{A'}
  = 0
\end{align*}
since  the one excitation state $\ket{\mathbf{j}} \neq \ket{\mathbf{0}}$ and so no $\ket{\mathbf{\ell }}$ which survive in both case.
Now for the one excitation product
\begin{equation*}
  \mathrm{Tr} _{A'} \ket{\mathbf{j}}\bra{\mathbf{k}}
  = \sum_{\boldsymbol{\ell}} \left( \mathbf{1 }_\mathbf{r}\otimes\bra{\boldsymbol{\ell}}_{A'} \right)\ket{\mathbf{j}}\bra{\mathbf{k}} \left( \mathbf{1 }_\mathbf{r}\otimes\ket{\boldsymbol{\ell}}_{A'} \right)
\end{equation*}
There are three possibilities, first $\mathbf{j} = \mathbf{k} = \mathbf{r}$
\begin{align*}
  \mathrm{Tr} _{A'} \ket{\mathbf{0}}\bra{\mathbf{j}}
   & =  \sum_{\boldsymbol{\ell}} \left( \mathbf{1 }_\mathbf{r} \bra{\boldsymbol{\ell}}_{A'} \right)\left(\ket{1}_\mathbf{r} \ket{\mathbf{0}}_{A'} \right) \left( \bra{{1}}_\mathbf{r}  \ket{\mathbf{0}}_{A'}\right)  \left( \mathbf{1 }_\mathbf{r} \ket{\boldsymbol{\ell}}_{A'} \right) \\
   & = \sum_{\boldsymbol{\ell}}  \ket{0}_\mathbf{r} \braket{\boldsymbol{\ell  }  |\mathbf{0} }_{A'} \bra{1}_\mathbf{r} \braket{\mathbf{0 } | \boldsymbol{\ell }}_{A'}
  = \ket{1}_\mathbf{r} \bra{1}_\mathbf{r}
\end{align*}
Next is $\mathbf{j } =\mathbf{k } \neq \mathbf{r}$
\begin{align*}
  \mathrm{Tr} _{A'} \ket{\mathbf{0}}\bra{\mathbf{j}}
   & =  \sum_{\boldsymbol{\ell}} \left( \mathbf{1 }_\mathbf{r} \bra{\boldsymbol{\ell}}_{A'} \right)\left(\ket{0}_\mathbf{r} \ket{\mathbf{j}}_{A'} \right) \left( \bra{{0}}_\mathbf{r}  \bra{\mathbf{k}}_{A'}\right)  \left( \mathbf{1 }_\mathbf{r} \ket{\boldsymbol{\ell}}_{A'} \right) \\
   & = \sum_{\boldsymbol{\ell}}  \ket{0}_\mathbf{r} \braket{\boldsymbol{\ell  }  |\mathbf{j} }_{A'} \bra{0}_\mathbf{r} \braket{\mathbf{k } | \boldsymbol{\ell }}_{A'}
  = \ket{0}_\mathbf{r} \bra{0}_\mathbf{r} \delta_{\mathbf{jk}}
\end{align*}
Finally, $\mathbf{j } \neq \mathbf{k } $, which is already covered here.
Putting them all together
\begin{gather*}
  \mathrm{Tr} _{A'} \ket{\mathbf{0}}\bra{\mathbf{0}}  = \ket{0}_\mathbf{r } \bra{0}_\mathbf{r} \\
  \mathrm{Tr} _{A'} \ket{\mathbf{0}}\bra{\mathbf{j}}  = \delta_{\mathbf{jr}} \ket{0}_\mathbf{r } \bra{1}_\mathbf{r} \\
  \mathrm{Tr} _{A'} \ket{\mathbf{j}}\bra{\mathbf{0}} = \delta_{\mathbf{jr}} \ket{1}_\mathbf{r } \bra{0}_\mathbf{r} \\
  \mathrm{Tr} _{A'} \ket{\mathbf{0}}\bra{\mathbf{j}} =\delta_{\mathbf{jk}} [(1-\delta_{\mathbf{jr}}) \ket{0}_\mathbf{r}\bra{0}_\mathbf{r}+\delta_{\mathbf{jr}} \ket{1}_\mathbf{r} \bra{1}_\mathbf{r}]
\end{gather*}
Therefore, the density operator
\begin{equation*}
  \rho(t) = \alpha^2   \ket{\mathbf{0}}\bra{\mathbf{0}} + \alpha \beta^* \sum_k f_k^*  \ket{\mathbf{0}}\bra{\mathbf{k}} + \alpha \beta \sum_j f_j \ket{\mathbf{j}}\bra{\mathbf{0}} + |\beta|^2 \sum_{j,k} f_j f_k^* \ket{\mathbf{j}}\bra{\mathbf{k}}
\end{equation*}
has the reduced state of
\begin{align*}
  \rho_\text{out}
   & = \alpha^2 \ket{0}_\mathbf{r } \bra{0}_\mathbf{r} + \alpha\beta^* \sum_k f_k^*  \delta_{kr}\ket{0}_\mathbf{r}\bra{1}_\mathbf{r} + \alpha \beta \sum_j f_j \delta_{jr}\ket{1}_\mathbf{r}\bra{0}_\mathbf{r} \\
   & \qquad + |\beta|^2 \sum_{j,k} f_j f_k^* \delta_{ {jk}} [(1-\delta_{ {jr}}) \ket{0}_\mathbf{r}\bra{0}_\mathbf{r}+\delta_{ {jr}} \ket{1}_\mathbf{r} \bra{1}_\mathbf{r}]                                     \\
   & = \alpha^2 \ket{0}_\mathbf{r } \bra{0}_\mathbf{r} + \alpha\beta^*  f_r^*\ket{0}_\mathbf{r}\bra{1}_\mathbf{r} + \alpha \beta f_r\ket{1}_\mathbf{r}\bra{0}_\mathbf{r}                                       \\
   & \qquad + |\beta|^2 \sum_{j \neq r} |f_j|^2 \ket{0}_\mathbf{r}\bra{0}_\mathbf{r} + |\beta|^2 |f_r|^2 \ket{1}_\mathbf{r} \bra{1}_\mathbf{r}                                                                 \\
  \rho_\text{out}
   & = \begin{bmatrix}
         \alpha^2 + |\beta|^2 \displaystyle\sum_{j \neq r} |f_j|^2 & \alpha\beta^*  f_r^* \\
         \alpha\beta  f_r                                          & |\beta|^2 |f_r|^2
       \end{bmatrix}
\end{align*}
The total normalized transition amplitude must equal to one,
\begin{equation*}
  \sum_j |f_j|^2 = \sum_j \braket{\mathbf{s}|U^\dagger|\mathbf{j} } \braket{\mathbf{j}|U|\mathbf{s}} = \braket{\mathbf{s}|U^\dagger U|\mathbf{s} } =1
\end{equation*}
so we can write the (0,0) element as
\begin{equation*}
  \rho_\text{out}
  = \begin{bmatrix}
    \alpha^2 + |\beta|^2 (1-|f_r|^2) & \alpha\beta^*  f_r^* \\
    \alpha\beta  f_r                 & |\beta|^2 |f_r|^2
  \end{bmatrix}
\end{equation*}

By defining output state
\begin{equation*}
  \ket{\psi_{\text{out}}(t)} = \frac{1}{\sqrt{P(t)}}\left(\cos\frac{\theta}{2} \ket{0} + e^{i\phi}\sin\frac{\theta}{2}f_r(t)\ket{1}\right) = \frac{1 }{\sqrt{P}} (\alpha \ket{0 }+\beta f_r \ket{1})
\end{equation*}
we can rewrite the reduced density operator
\begin{align*}
  \rho_\text{out}
   & =  P(t) \ket{\psi_{\text{out}}(t)} \bra{\psi_{\text{out}}(t)} + (1-P(t)) \ket{0} \bra{0} \\
   & =  \begin{bmatrix}
          \alpha^2         & \alpha\beta^*  f_r^* \\
          \alpha\beta  f_r & |\beta|^2 |f_r|^2
        \end{bmatrix}  +
  \begin{bmatrix}
    1-P(t) & 0 \\
    0      & 0
  \end{bmatrix}                                                                              \\
   & =  \begin{bmatrix}
          \alpha^2 +1 - P(t) & \alpha\beta^*  f_r^* \\
          \alpha\beta  f_r   & |\beta|^2 |f_r|^2
        \end{bmatrix}
\end{align*}
where it is understood $\ket{0}_\mathbf{r} = \ket{0}$.
Equating the (0,0)  element
\begin{align*}
  P & = 1 - |\beta|^2\bigl(1 - |f_r|^2\bigr) = 1 - |\beta|^2 + |\beta|^2|f_r|^2 \\
  P & = \alpha^2 + |\beta|^2|f_r|^2
\end{align*}

The fidelity is computed as
\begin{align*}
  F & =  \braket{\psi_\text{in} | \rho_{\text{out}}|\psi_\text{in} }                                            \\
    & = \begin{bmatrix}
          \alpha & \beta ^*
        \end{bmatrix}
  \begin{bmatrix}
    \alpha^2 + |\beta|^2 (1-|f_r|^2) & \alpha\beta^*  f_r^* \\
    \alpha\beta  f_r                 & |\beta|^2 |f_r|^2
  \end{bmatrix}
  \begin{bmatrix}
    \alpha \\  \beta
  \end{bmatrix}                                                                                               \\
    & = \begin{bmatrix}
          \alpha & \beta ^*
        \end{bmatrix}
  \begin{bmatrix}
    \alpha^3 + \alpha|\beta|^2 (1-|f_r|^2 + f_r ^*) \\
    \alpha^2\beta  f_r + |\beta|^2 |f_r|^2\beta
  \end{bmatrix}                                                               \\
    & =     \alpha^4 + \alpha^2|\beta|^2 (1-|f_r|^2 + f_r ^*) +    \alpha^2|\beta|^2  f_r + |\beta|^4 |f_r|^2   \\
  F & =     \alpha^4 + \alpha^2|\beta|^2 (1-|f_r|^2 ) + |\beta|^4 |f_r|^2 + 2\alpha^2|\beta|^2 \mathrm{Re}(f_r)
\end{align*}
The fidelity then averaged over the Bloch sphere
\begin{align*}
  \braket{F } = \frac{1 }{4\pi} \int_\Omega F \;d\Omega = \frac{1 }{2 } \int_0^\pi F(\theta) \sin \theta\;d\theta
\end{align*}
Recall that the total solid angle of a unit sphere is $4\pi$, thus the normalization constant.
Few helpful average value may be computated as follows
\begin{align*}
  \braket{\cos^4 \theta/2 }
   & = \frac{1 }{2 } \int_{0 }^{\pi} \cos^4 \frac{\theta }{2} \sin \theta\;d\theta
  = \frac{1 }{2 } \int_{0 }^{\pi} \left( \frac{1 +\cos \theta }{2} \right)^2  \sin \theta\;d\theta                                                                       \\
   & = \frac{1}{8}\left[\int_0^\pi \sin\theta\,d\theta + 2\int_0^\pi \cos\theta\sin\theta\,d\theta + \int_0^\pi \cos^2\theta\sin\theta\,d\theta\right]                   \\
  \braket{\sin^4 \theta/2 }
   & = \frac{1}{8} \left[ \cos \theta\big|_\pi^0 +2 \int_{-1 }^{1 }u\;du + \int_{-1 }^{1 }u^2\;du\right]  = \frac{1 }{8 }\left[ 2+0+\frac{2 }{3} \right]  = \frac{1 }{3}
\end{align*}
with $u = \cos\theta$ and $du = -\sin \theta \;d\theta$.
Other
\begin{align*}
  \braket{\sin^4 \theta/2 }
   & = \frac{1 }{2 } \int_{0 }^{\pi} \sin^4 \frac{\theta }{2} \sin \theta\;d\theta
  = \frac{1 }{2 } \int_{0 }^{\pi} \left( \frac{1 - \cos \theta }{2} \right)^2  \sin \theta\;d\theta                                                                     \\
   & = \frac{1}{8}\left[\int_0^\pi \sin\theta\,d\theta - 2\int_0^\pi \cos\theta\sin\theta\,d\theta + \int_0^\pi \cos^2\theta\sin\theta\,d\theta\right]                  \\
  \braket{\sin^4 \theta/2 }
   & = \frac{1}{8} \left[ \cos \theta\big|_\pi^0 + \int_{1 }^{-1 }u\;du + \int_{-1 }^{1 }u^2\;du\right]  = \frac{1 }{8 }\left[ 2+0+\frac{2 }{3} \right]  = \frac{1 }{3}
\end{align*}
and
\begin{align*}
  \braket{\cos ^2\theta/2\sin^2 \theta/2 }
   & =  \frac{1 }{2 }\int_{0 }^{\pi }\frac{1}{4 }(1+\cos \theta)(1-\cos \theta)\sin \theta\;d\theta          \\
   & =  \frac{1 }{8 }\int_{0 }^{\pi }(1-\cos^2 \theta)\sin \theta\;d\theta                                   \\
  \braket{\cos ^2\theta/2\sin^2 \theta/2 }
   & =  \frac{1 }{8 }\int_{-1 }^{1 }(1-u^2)\;du = \frac{1 }{8 }\left[ 2- \frac{2 }{3}\right]  = \frac{1 }{6} \\
\end{align*}
Then, with the fidelity as
\begin{equation*}
  F =     \cos^4 \frac{\theta}{2 } + \cos^2\frac{\theta}{2 }\sin ^2\frac{\theta}{2 } (1-|f_r|^2 ) + \sin ^4\frac{\theta}{2 } |f_r|^2 + 2\cos^2\frac{\theta}{2 }\sin ^2\frac{\theta}{2 } \mathrm{Re}(f_r)
\end{equation*}
The average is simply
\begin{equation*}
  \braket{F } = \frac{1 }{3 } + \frac{1 }{6} (1-|f_r|^2 ) + \frac{1 }{3 }|f_r|^2 + \frac{2 }{6 }\mathrm{Re }(f_r) = \frac{1 }{2 } + \frac{1 }{3} \mathrm{Re} (f_r) +\frac{1 }{6 }|f_r|^2
\end{equation*}

\subsubsection{Derivation: Purity of Bob's state.}
With the reduced state of
\begin{equation*}
  \rho_\text{out}
  = \begin{bmatrix}
    \alpha^2 + |\beta|^2 (1-|f_r|^2) & \alpha\beta^*  f_r^* \\
    \alpha\beta  f_r                 & |\beta|^2 |f_r|^2
  \end{bmatrix}
  = \begin{bmatrix}
    1- |\beta|^2 |f_r|^2 & \alpha\beta^*  f_r^* \\
    \alpha\beta  f_r     & |\beta|^2 |f_r|^2
  \end{bmatrix}
\end{equation*}
we can compute the purity as follows
\begin{align*}
  \mathrm{Tr} [\rho_\text{out}^2]
   & = \mathrm{Tr}
    \begin{bmatrix}
    1- |\beta|^2 |f_r|^2 & \alpha\beta^*  f_r^* \\
    \alpha\beta  f_r     & |\beta|^2 |f_r|^2
  \end{bmatrix}
  \begin{bmatrix}
    1- |\beta|^2 |f_r|^2 & \alpha\beta^*  f_r^* \\
    \alpha\beta  f_r     & |\beta|^2 |f_r|^2
  \end{bmatrix}                                                                  \\
   & =\mathrm{Tr} \begin{bmatrix}
                    (1-|\beta|^{2}|f_{r}|^{2})^2 + |\alpha|^{2}|\beta|^{2}|f_{r}|^{2} & \alpha\beta^{*}f_{r}^{*} \\
                    \alpha\beta f_{r}                                                 &
                    |\alpha|^{2}|\beta|^{2}|f_{r}|^{2}+|\beta|^{4}|f_{r}|^{4}
                  \end{bmatrix}\\
  & = 1+ |\beta|^{4}|f_{r}|^{4} - 2|\beta|^{2}|f_{r}|^{2} +   2|\alpha|^{2}|\beta|^{2}|f_{r}|^{2} + |\beta|^{4}|f_{r}|^{4}\\
  & = 1+ 2|\beta|^{4}|f_{r}|^{4} - 2|\beta|^{2}|f_{r}|^{2} +   2|\beta|^{2}|f_{r}|^{2} - 2|\beta|^{4}|f_{r}|^{2} \\
  & = 1+ 2|\beta|^{4}|f_{r}|^{2}(|f_r| ^2 -1)  \\
\end{align*}


\subsubsection{Derivation: Classical fidelity.}
The classical density operator can be written as ensemble of pure state
\begin{equation*}
  \rho_\text{cl} (\phi) = \sum_{k} p_k \ket{\phi_k} \bra{\phi_k}
\end{equation*}
It is classical–quantum hybrid object, not a quantum state evolving under Hamiltonian dynamics.
The probability $p_k$ is a simple Born's rule
\begin{align*}
  p_k & =  \braket{\psi  | E_k | \psi} = \mathrm{Tr} (E_k \rho_\psi)                                                                                                                                                        \\
      & =\frac{1 }{2} \mathrm{Tr}\big[(a_k \mathbf{1} +  \mathbf{b}_k\cdot \boldsymbol{\sigma})(\mathbf{1} + \mathbf{n}\cdot\boldsymbol{\sigma})\big],                                                                      \\
      & = \frac{1  }{2 }\mathrm{Tr}\big[a_k \mathbf{1} + a_k \mathbf{n }\cdot\boldsymbol{\sigma} + (\mathbf{b}_k\cdot\boldsymbol{\sigma}) + (\mathbf{b}_k\cdot\boldsymbol{\sigma})(\mathbf{n}\cdot\boldsymbol{\sigma})\big]
\end{align*}
We write $E_k$ can in terms of Pauli matrices, along with the identity, since they span the $2 \times 2$ matrix space.
Also $\ket{\psi}$ here denote the actual quantum state.
The first term is trivial, the trace of identity is the dimension.
As for second term, we note that the trace of Pauli matrix is zero
\begin{equation*}
  \mathrm{Tr} a_k \mathbf{n }\cdot \boldsymbol{\sigma} = a_k \mathrm{Tr} \sum_i n_i \sigma_i = a_k \sum_{i}  n_i  \mathrm{Tr} \sigma_i  = 0
\end{equation*}
and the same for the third term, we use the identity $    (\mathbf{a} \cdot \boldsymbol{\sigma})(\mathbf{b} \cdot \boldsymbol{\sigma}) = \mathbf{a} \cdot \mathbf{b} + i(\mathbf{a} \times \mathbf{b}) \cdot \boldsymbol{\sigma}$
\begin{align*}
  \mathrm{Tr} \varsigma & =  \mathrm{Tr} (\mathbf{n} \cdot \mathbf{b}_k \mathbf{1}) + \mathrm{Tr} i(\mathbf{n} \times \mathbf{b}_k) \cdot \boldsymbol{\sigma} \\
                        & =  2\mathbf{n} \cdot \mathbf{b}_k  +i \sum_i (b_k \times n)_i \,\mathrm{Tr}(\sigma_i)
  =  2\mathbf{n} \cdot \mathbf{b}_k
\end{align*}
So $p_k = a_k + \mathbf{n }\cdot \mathbf{b}_k$.
Substituting this
\begin{equation*}
  \rho_\text{cl} = \frac{1 }{2 }  \sum_k \left[ \left( a_k  \mathbf{1}+ \mathbf{n }\cdot \mathbf{b}_k \right)    (\mathbf{1 }+ \mathbf{r}_k \cdot \boldsymbol{\sigma})\right]
\end{equation*}
Note that $a_k$ and $\mathbf{b}_k$ belong to the element of $E_k$, while $\mathbf{r}_k$ is the Bloch vector of prepared output state.
Once the protocol is determine, $\mathbf{r}_k$ is fixed

Next is to compute the fidelity
\begin{align*}
  F & =  \braket{\psi| \rho_\text{cl} | \phi} = \mathrm{Tr} (\rho_{\phi}\rho_\text{cl} )                                                                                                                                                                                                         \\
    & = \frac{1 }{4 } \mathrm{Tr} \sum_k \left( \mathbf{1} + \mathbf{n }\cdot \boldsymbol{\sigma} \right) \left( a_k  \mathbf{1}+ \mathbf{n }\cdot \mathbf{b}_k \right)    (\mathbf{1 }+ \mathbf{r}_k \cdot \boldsymbol{\sigma})                                                                 \\
    & = \frac{1 }{4 } \sum_k \left( a_k  \mathbf{1}+ \mathbf{n }\cdot \mathbf{b}_k \right) \mathrm{Tr}  \left[\left( \mathbf{1} + \mathbf{n }\cdot \boldsymbol{\sigma} \right) (\mathbf{1 }+ \mathbf{r}_k \cdot \boldsymbol{\sigma})\right]                                                      \\
    & = \frac{1 }{4 } \sum_k \left( a_k  \mathbf{1}+ \mathbf{n }\cdot \mathbf{b}_k \right) \mathrm{Tr}  \left[\mathbf{1} + \mathbf{r}_k \cdot \boldsymbol{\sigma} + \mathbf{n }\cdot \boldsymbol{\sigma} + (\mathbf{n }\cdot \boldsymbol{\sigma})(\mathbf{r}_k \cdot \boldsymbol{\sigma})\right] \\
    & = \frac{1 }{4 } \sum_k \left( a_k  \mathbf{1}+ \mathbf{n }\cdot \mathbf{b}_k \right) \left( 2 + 2 \mathbf{n }\cdot \mathbf{r}_k\right)                                                                                                                                                     \\
  F & = \frac{1 }{2 }  \sum_k\left[ a_k + a_k(\mathbf{n }\cdot \mathbf{r }_k) + (\mathbf n\cdot\mathbf b_k) + (\mathbf n\cdot\mathbf b_k)(\mathbf{n }\cdot \mathbf{r }_k)\right]
\end{align*}
The averaged value is then
\begin{equation*}
  \braket{F } = \frac{1 }{8 \pi } \sum_{k} \int d\Omega\;\left[ a_k + a_k(\mathbf{n }\cdot \mathbf{r }_k) + (\mathbf n\cdot\mathbf b_k) + (\mathbf n\cdot\mathbf b_k)(\mathbf{n }\cdot \mathbf{r }_k)\right]
\end{equation*}
The averaged fidelity compute average over all possible unknown input states, which parameterized by $\mathbf{n} = (\sin \theta \cos \phi,\sin \theta \sin\phi,\cos \theta)$.
This means that $d\Omega$ acts only on $\mathbf{n}$, not $a_k$ nor $\mathbf{r}_k$.
We evaluate term by term starting from the first
\begin{equation*}
  \braket{F }_1 = \frac{1 }{8 \pi } \sum_{k} \int d\Omega\; a_k = \frac{1 }{2 }\sum_{k } a_k
\end{equation*}
The second
\begin{align*}
  \braket{F }_2 & =  \frac{1 }{8 \pi } \sum_{k} \int d\Omega\; a_k(\mathbf{n }\cdot \mathbf{r }_k)
  = \frac{1 }{8 \pi } \sum_{k,l} \int d\Omega\; a_k n_l r_{rl}                                     \\
                & =  \frac{1 }{8 \pi } \sum_{k,l} a_k r_{rl} \int d\Omega\; n_l = 0
\end{align*}
The result is due to the integral being zero.
This can be understood as that for every $\mathbf{n}$, there exist $-\mathbf{n}$ such that the average over the solid angle is zero.
Same reasoning applies to the third term.
For the fourth
\begin{align*}
  \braket{F }_4
   & = \frac{1 }{8 \pi } \sum_{k} \int d\Omega\;  (\mathbf n\cdot\mathbf b_k)(\mathbf{n }\cdot \mathbf{r }_k) \\
   & = \frac{1 }{8 \pi } \sum_{k,l,m} \int d\Omega\;  n_lb_{rl}n_{m}r_{km}                                    \\
   & = \frac{1 }{8 \pi } \sum_{k,l,m} b_{rl}r_{km}\int d\Omega\;  n_ln_{m}                                    \\
\end{align*}
To evaluate the integral, we use the isotropic identity
\begin{equation*}
  n_x^2 + n_y^2 + n_z^2 = 1\qquad\langle n_x^2\rangle = \langle n_y^2\rangle = \langle n_z^2\rangle = \frac{1}{3}\qquad
  \frac{1 }{4\pi}\int d\Omega\;  n_j n_{j} = \frac{1 }{3} \delta_{ij}
\end{equation*}
Then
\begin{equation*}
  \braket{F }_4
  = \frac{1 }{2} \sum_{k,l,m} b_{rl}r_{km}\frac{1 }{3 } \delta_{lm} =  \frac{1 }{6 }\sum_{k } \mathbf{b }_k \cdot \mathbf{r}_k
\end{equation*}
and we obtain the average fidelity
\begin{equation*}
  \braket{F } = \frac{1 }{2 } \sum_{k}a_k + \frac{1 }{6 }\sum_{k } \mathbf{b }_k \cdot \mathbf{r}_k
\end{equation*}
As a unit vector, $\mathbf{r}_k$ and $\mathbf{b}_k$ has magnitude of unitary.
At maximum therefore, the value of average fidelity is $\braket{F } = 1/2+1/6=2/3$, lower than quantum fidelity.


\subsection{Ergotropy}
Suppose the system is in a state $\rho$ with average energy of $U(\rho) = \mathrm{Tr}  (\rho H)$.
Ergotropy is the maximum amount of work that can be extracted from a quantum system by a cyclic unitary process
\begin{equation*}
  \mathcal{E}(\rho) = \mathrm{Tr}[\rho \mathcal{H}] - \min_{U} \mathrm{Tr}[U_t \rho U_t^\dagger \mathcal{H}] 
\end{equation*}
The passive energy $U_p =  \min_{U} \mathrm{Tr}[U_t \rho U_t^\dagger \mathcal{H}] $ is the minimum energy reachable by unitary transformations.
Ergotropy can be stored in a system in two ways, either through a population inversion and/or in coherences in the energy eigenbasis

Recall that in statistical mechanics, heat corresponds to energy exchange between system and environment without altering its Hamiltonian, which then changes the population; while work represent changes in the energy spectrum while the occupations remain unchanged
\begin{equation*}
  dE = \sum_{s } (\epsilon_s\;dn_s+n_s\;\epsilon_s) = dQ+dW
\end{equation*}
Similarly, heat in quantum mechanics interpreted as changing occupation probabilities, while work changes the energy level 
\begin{equation*}
  dE = \mathrm{Tr} (d\rho\;H) + \mathrm{Tr} (\rho\;dH)= dQ+dW
\end{equation*}
In unitary transformations, the system is isolated from heat bath, $dQ = 0$, and all energy changes is due to work.
Both the state and Hamiltonian do change over time, it just that the work contribution is due to changes in Hamiltonian $dH$.
All this simply describe the protocol to extract energy from the system.

By the  definition, ergotropy is the maximum extractable work, as such arbitrary energy extraction does not yield the ergotropy value.
The protocol is as follows.
\begin{enumerate}
  \item Start with system at $\rho$ and apply external control (magnetic field, laser pulse, gate voltages, etc.), making the Hamiltonian time-dependent $H \to H(t)$, thus evolving the state $\rho \to U \rho U ^\dagger$.
  \item Choose the external control such that the generated unitary brings the system to the passive state $    \rho(T) = \rho_p$, while the Hamiltonian is returned to its original form $    H(T) = H(0) = H$.
  \item At this point, no further unitary operation can lower the energy. The system is work exhausted, and the energy extracted is the ergotropy
\end{enumerate}

\subsubsection{Pure and mixed state.}
Suppose we are given the pure and mixed state of spin chain in the form 
\begin{gather*}
  \ket{\psi(0)} = \left(\alpha\ket{0}_1 + \beta\ket{1}_1\right)\bigotimes_{i=2}^{N}\ket{0}_i  = \alpha\ket{\mathbf{0}} + \beta\ket{\mathbf{1}}\\
\rho(0) = \left(q\ket{1}\bra{1}_1 + (1-q)\ket{0}\bra{0}_1\right)\bigotimes_{i=2}^{N}\ket{0}\bra{0}_i = q\ket{\mathbf{1}}\bra{\mathbf{1}} + (1-q)\ket{\mathbf{0}}\bra{\mathbf{0}}
\end{gather*}
and we are tasked to determine the ergotropy of the same state.
The first lines written in such way to demonstrate that the first site of spin chain is initially pure and mixed respectively.
First we set the ground state as zero, that is $H \ket{\mathbf{0}} = E_0 \ket{\mathbf{0}} = 0$.
This is the passive state for the pure state, while the initial energy is 
\begin{equation*}
  E_i = |\alpha|^2 \bra{\mathbf{0}} H \ket{\mathbf{0}} + |\beta|^2 \bra{\mathbf{1}} H \ket{\mathbf{1} }= |\beta|^2 \bra{\mathbf{1}} H \ket{\mathbf{1} }
\end{equation*}
This is also the ergotropy.
The initial energy of the mixed state can be computed in similar manner 
\begin{equation*}
  E_i = \mathrm{Tr} (\rho(0)H) = q\mathrm{Tr} (\ket{\mathbf{1}}\bra{\mathbf{1}}H) + (1-q) \mathrm{Tr} (\ket{\mathbf{0}}\bra{\mathbf{0}}H) = q \braket{\mathbf{1 } | H| \mathbf{1}}
\end{equation*}
For passive state, however, things differ.
The ergotropy is stored via population inversion, so the passive state has the following form 
\begin{equation*}
  \rho(0) =    (1-q)\ket{\mathbf{1}}\bra{\mathbf{1}} + q\ket{\mathbf{0}}\bra{\mathbf{0}}
\end{equation*}
Then the passive energy is simply $E_p = (1-q) \braket{\mathbf{1 } | H| \mathbf{1}}$ and the ergotropy  is $\mathcal{E} =(2q-1) \braket{\mathbf{1 } | H| \mathbf{1}} $.
Suppose both case has the same initial ergotropy, then $q = (1+|\beta^2|)/2$.

Let us then write the density matrix in the following form 
\begin{equation*}
  \rho_\text{out}
  = \begin{bmatrix}
    \alpha^2 + |\beta|^2 (1-|f_r|^2) & \alpha\beta^*  f_r^* \\
    \alpha\beta  f_r                 & |\beta|^2 |f_r|^2
  \end{bmatrix}
\end{equation*}
Of course for the case of mixed state, the off diagonal equal to zero.
Then the product of the density state with the local Hamiltonian $H = -B\sigma_z$ is 
\begin{equation*}
  \rho H = \begin{bmatrix}
    0&-B\rho_{01}\\
    0&-B\rho_{11}
  \end{bmatrix}
\end{equation*}
Hence, the initial energy is $E_i = -B |\beta|^2 |f_n|$ for both the pure and mixed state.
Next we consider the passive energy of the pure state, which is obtained by putting the largest eigenvalue of $\rho$ into the lowest energy level.
The eigenvalue of $\rho$ as $2 \times 2$ matrix has the following form 
\begin{equation*}
  \lambda_{\pm} = \frac{1 }{2 } \sqrt{1 \pm \sqrt{1-4 \det \rho}}
\end{equation*}
whereas the determinant is 
\begin{align*}
  \det \rho & = [\alpha^2 + |\beta|^2 (1-|f_n|^2)]|\beta|^2|f_n|^2 - |\alpha|^2|\beta^2||f_n|^2\\ 
  \det \rho & = |\beta|^4|f_n|^2 (1-|f_n|^2) 
\end{align*}
So 
\begin{equation*}
  \lambda_{\pm} = \frac{1 }{2 } \left[ 1 \pm \sqrt{1+4 |\beta|^4|f_n|^2 (|f_n|^2-1)} \right] 
\end{equation*}
and we have the passive energy as $E_p = -B\lambda_-$.
Using the initial and passive energy, we find that the ergotropy for coherent state is 
\begin{equation*}
  \mathcal{E}_\text{coh} (t) = 2|f_n(t)|^2\sin^2\left(\frac{\theta}{2}\right)-1+\sqrt{1+4\sin^4\left(\frac{\theta}{2}\right)|f_n(t)|^2\left(|f_n(t)|^2-1\right)}
\end{equation*}
after substituting $\alpha$ and $\beta$ parameter along with choosing $B=2$.
As for mixed state, the ergotropy is stored in population inversion.
Assuming that the lower state population $|\beta|^2|f_n|^2>0.5$, the passive energy after population inversion is $E_p = -B[\alpha^2 + |\beta|^2 (1-|f_r|^2) ]$ and the ergotropy of 
\begin{equation*}
  \mathcal{E}_\text{mix} (t) = [2+2\sin (\theta/2)^2]|f_n|^2  - 2
\end{equation*}

\end{document}